% Source: https://artofproblemsolving.com/community/c3h1547785_2004_mathcounts_nats
% Stable identity: aops-topic-1547785-q26
Since $a^2b^3c=2^2\cdot3^3\cdot5^7$, the numbers $a$, $b$, $c$ have no prime factors other than $2,3,5$. Write the exponent of a prime in $a$, $b$ as $\alpha$, $\beta$; then $2\alpha+3\beta$ is at most the exponent of that prime in the product, and the rest goes into $c$.
\begin{itemize}
\item Prime $2$: $2\alpha+3\beta\le2$, so $\beta=0$, $\alpha\le1$, and $\alpha+\beta\le1$.
\item Prime $3$: $2\alpha+3\beta\le3$, so $(\alpha,\beta)$ is $(0,0)$, $(1,0)$ or $(0,1)$, and $\alpha+\beta\le1$.
\item Prime $5$: $2\alpha+3\beta\le7$; the largest value of $\alpha+\beta$ is $3$, from $(2,1)$ or $(3,0)$.
\end{itemize}
Hence $ab\le2\cdot3\cdot5^3=750$, and since $c\ge1$,
\[
ab-c\le750-1=749 .
\]
Equality is possible: $a=2\cdot5^2=50$, $b=3\cdot5=15$, $c=1$ gives $a^2b^3=2^2\cdot5^4\cdot3^3\cdot5^3=2^2\cdot3^3\cdot5^7$ exactly. The greatest value is $50\cdot15-1=\boxed{749}$.
